\section{Numerical method}
\label{sec:numerics}

\NHFOAM\ discretises \eqref{eq:cont}--\eqref{eq:mom} with OpenFOAM's
cell-centred finite-volume method on the moving $\sigma$-mesh, and couples
pressure and velocity with a PIMPLE (PISO-type) algorithm \citep{issa1986}.
What is specific to the model is where the pressure is stored, how the free
surface enters the pressure equation, and how the time step is organised.

\subsection{Pressure on the layer interfaces}
\label{sec:interface}

Two vertical arrangements of the pressure are available
(\code{pressureLocation}). The default, \code{interface}, follows NHWAVE and
SWASH \citep{stelling2003,zijlema2005,ma2012}: $\prgh$ is stored on the layer
interfaces ($\sigma$-levels), Fig.~\ref{fig:interface}. There is one unknown on
the bottom face of every cell (the bed and the internal interfaces) and one
surface unknown per column, $p_s = g\,\eta^{n+1}$, so that $q=0$ holds exactly
at the free surface \eqref{eq:dynamic}.

\begin{figure}[t]
\centering
\begin{tikzpicture}[x=1cm,y=1cm,>=Stealth]
  % two columns, three layers, sloping bottom and surface
  \def\xa{0}\def\xb{3.6}\def\xc{7.2}
  \coordinate (b0) at (\xa,0.0);  \coordinate (b1) at (\xb,0.45); \coordinate (b2) at (\xc,0.9);
  \coordinate (s0) at (\xa,4.8);  \coordinate (s1) at (\xb,5.0);  \coordinate (s2) at (\xc,4.9);
  \foreach \k in {0,1,2,3} {
    \draw[thick] ($(b0)!\k/3!(s0)$) -- ($(b1)!\k/3!(s1)$) -- ($(b2)!\k/3!(s2)$);
  }
  \draw[thick] (b0) -- (s0); \draw[thick] (b1) -- (s1); \draw[thick] (b2) -- (s2);
  % interface CV of the middle interface of the left column (shaded)
  \fill[blue!12] ($(b0)!0.5/3!(s0)$) -- ($(b1)!0.5/3!(s1)$) -- ($(b1)!1.5/3!(s1)$) -- ($(b0)!1.5/3!(s0)$) -- cycle;
  \foreach \k in {0,1,2,3} {
    \draw[thick] ($(b0)!\k/3!(s0)$) -- ($(b1)!\k/3!(s1)$);
  }
  % interface unknowns (circles) at face centres of layer faces, surface unknowns (squares)
  \foreach \k in {0,1,2} {
    \fill[blue!70!black] ($($(b0)!\k/3!(s0)$)!0.5!($(b1)!\k/3!(s1)$)$) circle (2.6pt);
    \fill[blue!70!black] ($($(b1)!\k/3!(s1)$)!0.5!($(b2)!\k/3!(s2)$)$) circle (2.6pt);
  }
  \fill[red!70!black] ($($(s0)!0.5!(s1)$)+(-0.09,-0.09)$) rectangle +(0.18,0.18);
  \fill[red!70!black] ($($(s1)!0.5!(s2)$)+(-0.09,-0.09)$) rectangle +(0.18,0.18);
  % cell centres (crosses) of the left column
  \foreach \k in {0.5,1.5,2.5} {
    \node at ($($(b0)!\k/3!(s0)$)!0.5!($(b1)!\k/3!(s1)$)$) {$\times$};
  }
  % side flux, split half/half
  \draw[->,thick,orange!80!black] ($($(b1)!1.0/3!(s1)$)!0.5!($(b1)!2.0/3!(s1)$)+(-0.45,0)$) -- +(0.9,0)
     node[pos=0.8,above,font=\footnotesize,text=orange!80!black]{$\phi_{f}$};
  % mid-plane vertical flux
  \draw[->,thick,green!50!black] ($($(b0)!1.5/3!(s0)$)!0.5!($(b1)!1.5/3!(s1)$)+(0,-0.42)$) -- +(0,0.84);
  \node[font=\footnotesize,green!40!black,anchor=west] at ($($(b0)!1.5/3!(s0)$)!0.5!($(b1)!1.5/3!(s1)$)+(0.08,0.32)$) {$F_{\mathrm{mid}}$};
  % labels
  \node[font=\footnotesize,anchor=east] at ($(s0)+(-0.1,0)$) {$\eta$};
  \node[font=\footnotesize,anchor=east] at ($(b0)+(-0.1,0)$) {$z_b$};
  \node[font=\footnotesize,anchor=west,align=left] at (7.6,4.6)
     {\textcolor{red!70!black}{$\blacksquare$} surface unknown $p_s=g\eta^{n+1}$};
  \node[font=\footnotesize,anchor=west,align=left] at (7.6,4.0)
     {\textcolor{blue!70!black}{$\bullet$} interface unknown (bottom face\\ \quad of each cell)};
  \node[font=\footnotesize,anchor=west,align=left] at (7.6,3.15)
     {$\times$ cell centre: mean of its\\ \quad two interfaces};
  \node[font=\footnotesize,anchor=west,align=left] at (7.6,2.3)
     {\colorbox{blue!12}{\phantom{x}} interface control volume};
\end{tikzpicture}
\caption{Pressure on the layer interfaces: two columns of three layers. Each
interface carries one unknown; the shaded volume, half a cell below and half a
cell above an interface, is where continuity is imposed. A side flux $\phi_f$
is shared equally by the interfaces below and above it, and the flux
$F_{\mathrm{mid}}$ through the sloping mid-plane of a cell connects its two
interfaces.}
\label{fig:interface}
\end{figure}

\paragraph{Gradient.} Inside a cell, the vertical gradient is compact,
$(p_{\mathrm{top}}-p_{\mathrm{bot}})/\Delta z$. The value at the cell centre is
the mean of the two interfaces, and the horizontal gradient is the mean of the
two interface gradients, as in NHWAVE. In practice the momentum equation takes
the Gauss gradient with the interface values on the layer faces and the
interpolated centre values on the side faces.

\paragraph{Continuity.} Continuity is imposed on interface-centred control
volumes, half a cell below and half a cell above each interface, as in
NHWAVE's Poisson equation. With $\bm{H}/a_P$ the velocity predicted without the
pressure gradient and $r_A = 1/a_P$ the inverse of the momentum diagonal, the
flux through a side face $f$ between cells $P$ and $N$ is
\begin{equation}
  \phi_f = \phi^{*}_f - (r_A)_f\,|\bm{S}_f|\,\delta_f\left(\bar p_N - \bar p_P\right),
  \qquad \bar p = \tfrac12\left(p_{\mathrm{bot}} + p_{\mathrm{top}}\right),
  \label{eq:sideflux}
\end{equation}
so it couples the four interfaces of its two cells. Half of $\phi_f$ goes to
the control volume of the interface below and half to the one above. The flux
through the sloping mid-plane of cell $P$, of area vector
$\bm{S}_{\mathrm{mid}} = (S_x,S_y,A)$ with $A$ the horizontal area of the
column, connects the cell's two interfaces. It is split as OpenFOAM's
over-relaxed non-orthogonal Laplacian \citep{jasak1996}:
\begin{equation}
  F_{\mathrm{mid}} = \frac{\bm{H}}{a_P}\cdot\bm{S}_{\mathrm{mid}}
  - r_A\,\frac{A^2 + |\bm{S}_h|^2}{V}\,(p_{\mathrm{top}} - p_{\mathrm{bot}})
  - r_A\left[S_x\frac{\partial p}{\partial x} + S_y\frac{\partial p}{\partial y}
  - \frac{|\bm{S}_h|^2}{A}\frac{\partial p}{\partial z}\right]^{\mathrm{lagged}},
  \label{eq:midflux}
\end{equation}
where $|\bm{S}_h|^2 = S_x^2+S_y^2$ and $V$ is the cell volume. The implicit
coefficient carries the full $|\bm{S}_{\mathrm{mid}}|^2 = A^2 + |\bm{S}_h|^2$,
so the lagged bracket holds no vertical pressure difference, and the iteration
over the lagged terms converges however steep the layers are. With $A^2$ alone,
the pressure correctors diverged under the bore fronts of the breaking case
(Section~\ref{sec:breaking}).

\paragraph{Free surface.}
\label{sec:freesurface}
The control volume of the surface unknown is the upper half of the top cell.
Its outflow is the absolute flux through the surface face,
$\phi_s = A_h\,\partial\eta/\partial t$, discretised as
$A_h(\eta^{n+1}-\hat\eta)/a$ (Section~\ref{sec:time}). With
$\eta^{n+1}=p_s/g$ and the relaxation zones \eqref{eq:relax}, the equation of the
surface unknown of each column is
\begin{equation}
  \frac{A_h}{a}\left[(1 + a\sigma_r)\frac{p_s}{g} - \hat\eta - a\sigma_r\eta_t\right]
  = \tfrac12\sum_{\text{sides of the top cell}}\phi_f + F_{\mathrm{mid}} .
  \label{eq:surface}
\end{equation}
Each pressure solve therefore delivers $q$ and $\eta^{n+1}$ together. There is
no gravity-wave Courant limit, the matrix stays symmetric, and no pressure
reference is needed. After the solve, the side fluxes follow from
\eqref{eq:sideflux}, the fluxes through the layer faces are integrated upward
from the bed so that every cell conserves volume exactly, and
$\eta^{n+1}$ follows from the flux through the surface. In closed basins the
water volume is conserved to round-off.

\paragraph{Matrix and preconditioner.} The matrix of the interface unknowns is
symmetric positive definite, but it is not diagonally dominant where the
layers are thicker than the horizontal spacing, and incomplete Cholesky (DIC)
can then break down. \NHFOAM\ solves it with preconditioned conjugate gradients
and its own \code{sigmaColumn} preconditioner, which solves the tridiagonal
block of every column exactly (line relaxation in the vertical). It handles
thin layers and the strong vertical coupling of breaking columns, and falls
back to DIC (DILU for asymmetric matrices) on ordinary OpenFOAM matrices, so a
case can name it in either layout. Three kinds of terms would need more
couplings; they are lagged and updated in every non-orthogonal corrector:
\begin{enumerate}[nosep]
  \item the non-orthogonal correction of the side faces;
  \item the along-layer part of the mid-plane flux, the bracket in
        \eqref{eq:midflux};
  \item on processor boundaries, the couplings to the cell above the remote
        cell. The face is written in a compact form that bounds the exact one,
        and the difference is lagged.
\end{enumerate}

\subsection{Cell-centred pressure}
\label{sec:cellp}

With \code{pressureLocation cell}, $\prgh$ is OpenFOAM's usual cell-centred
field, and the kinematic and dynamic conditions combine into a Robin condition
on the surface patch, implemented through OpenFOAM's \code{mixed} patch field:
\begin{equation}
  p_s = f\,p_{\mathrm{ref}} + (1-f)\,p_c,\quad
  p_{\mathrm{ref}} = \frac{g\left(\hat\eta + a\sigma_r\eta_t + a\,\phi^{*}_s/A_h\right)}{1 + a\sigma_r},\quad
  f = \frac{1 + a\sigma_r}{1 + a\sigma_r + \beta_s},
  \label{eq:robin}
\end{equation}
with $\beta_s = g\,a\,(r_A)_f\,|\bm{S}_f|\,\delta_f/A_h$ and $p_c$ the value in the
top cell. This layout needs many more layers for the same dispersion
(Section~\ref{sec:dispersion}) and is kept for reference.

\subsection{Breaking columns}
\label{sec:numhfa}

In breaking columns the vertical coupling of the pressure is multiplied by a
large factor (\code{hydrostaticFactor}, default 1000). In the interface layout
this is the coefficient of $p_{\mathrm{top}}-p_{\mathrm{bot}}$ in
\eqref{eq:midflux}, up to the surface unknown; in the cell layout it is an
extra coupling across the layer faces plus a large mobility of the surface
face. $\prgh$ then becomes uniform in the vertical ($q=0$, since the surface
value is $g\eta$), and the fluxes through the layer faces follow from
continuity. The factor acts on the vertical pressure difference only: scaling
the face mobility itself would also scale the non-orthogonal correction on the
sloping layer faces, i.e.\ the horizontal pressure gradient under a steep
front, which drives spurious vertical fluxes.

The breaking flags, ages and eddy viscosities are set at the start of each
time step from $\partial\eta/\partial t = (\eta^n-\eta^{n-1})/\Delta t_0$.
The flags spread over the neighbouring columns in repeated sweeps until
nothing changes; on a processor they reach only the local columns. The
eddy viscosity of a face is the linear interpolate of the column values
(zero on walls, the free surface and the bottom, and on the closed faces of
dry columns), and the diffusion term is implicit in $\bm{u}$
(\code{fvm::laplacian}), so it needs no time-step limit.

\subsection{Time integration}
\label{sec:time}

Time integration follows the \code{ddt(U)} scheme: Euler, or the second-order
backward differentiation formula (BDF2) with variable time-step coefficients
for both $\bm{u}$ and $\eta$. With $r = \Delta t/\Delta t_0$, where $\Delta t_0$ is
the previous step,
\begin{equation}
  \begin{aligned}
  \left.\frac{\partial\eta}{\partial t}\right|^{n+1} &\approx \frac{\eta^{n+1}-\hat\eta}{a},\qquad
  a = \frac{\Delta t}{c},\qquad
  \hat\eta = \frac{c_0\,\eta^{n} - c_{00}\,\eta^{n-1}}{c},\\
  c &= 1 + \frac{r}{1+r},\qquad
  c_{00} = \frac{r^2}{1+r},\qquad
  c_0 = c + c_{00}.
  \end{aligned}
  \label{eq:bdf2}
\end{equation}
The first step after the start (or a restart) is Euler. OpenFOAM's Rhie--Chow
type time-derivative flux correction (\code{ddtCorr}) is used with its
coefficient set to one (\code{backward 1} in \code{ddtSchemes}); with the
default limiter, standing waves lost 3.9\% of their amplitude over 4.5
periods instead of 0.09\%.

\paragraph{Algorithm per time step.}
\begin{enumerate}[nosep]
  \item Shift the time levels of $\eta$ and of the face velocity $\bm{u}_f$; set the
        targets of the relaxation zones and the breaking flags.
  \item Move the mesh to the extrapolated surface
        $\eta^{*} = \eta^n + r\,(\eta^n - \eta^{n-1})$.
  \item Rebuild the flux on the moved faces from the face velocity and make
        it relative to the mesh motion. In the first outer corrector the face
        velocity is extrapolated (\code{fluxExtrapolation}, default on with
        BDF2):
        \begin{equation}
          \phi^{*} = \bm{S}_f\cdot\left[(1+r)\,\bm{u}_f^{n} - r\,\bm{u}_f^{n-1}\right] - \phi_{\mathrm{mesh}} .
          \label{eq:fluxext}
        \end{equation}
  \item Momentum predictor, with the pressure gradient of the last corrector.
  \item Pressure correctors (\code{nCorrectors}, 2--3). Each one assembles and
        solves the interface matrix with the surface rows \eqref{eq:surface}
        (or, in the cell layout, sets the Robin coefficients \eqref{eq:robin}
        and solves for $\prgh$), with \code{nNonOrthogonalCorrectors} updates
        of the lagged terms, then updates the flux, $\eta^{n+1}$, $\bm{u}$ and
        $\bm{u}_f$.
  \item Optional further outer correctors repeat steps 3--5 with the latest
        flux (and, with \code{moveMeshOuterCorrectors}, the latest $\eta$).
\end{enumerate}

\paragraph{Order of accuracy in time.} With one outer corrector, the momentum
equation of step $n+1$ is linearised with a convecting flux of level $n$. That
lag is an error of order $\Delta t$ in every step, so the scheme is only first
order in time for nonlinear waves, although BDF2 is second order. In the
solitary wave of Section~\ref{sec:ma-solitary}, halving the time step halved
the error. The extrapolation \eqref{eq:fluxext} removes the lag at no cost.
Before it was introduced, a second outer corrector gave the same accuracy at
1.6 times the cost (Section~\ref{sec:ma-solitary}).

\subsection{Convection}
\label{sec:convection}

The convection term uses OpenFOAM's schemes. Three are used in this report:
\begin{itemize}[nosep]
  \item \code{Gauss limitedLinearV 1}, a TVD scheme. It is needed in breaking
        cases, where bores form, and it was the default until this report;
  \item \code{Gauss filteredLinear2V 0.2 0}, central differencing with a filter
        that only acts on grid-scale oscillations. It damps waves as little as
        central differencing but stayed stable in long nonlinear runs
        (Section~\ref{sec:ma}); it is the recommended scheme for non-breaking
        waves;
  \item \code{Gauss linear}, pure central differencing. It damps least, but
        grid-scale noise grows in long nonlinear runs and behind bores.
\end{itemize}
Upwind-biased schemes (\code{linearUpwind}, \code{LUST}) damp waves of moderate
steepness noticeably (Section~\ref{sec:convtests}).
